\documentclass[12pt, oneside]{book} 
\pagestyle{plain}

% The \usepackage{} command will import predefined fonts, symbols, environments, etc.  For example, the ams packages below come from the American Mathematical Society and include all kinds of useful math symbols like integrals
\usepackage{amscd}
\usepackage{amsmath}
\usepackage{amssymb}
\usepackage{amsthm}
\usepackage{verbatim}
\usepackage[utf8]{inputenc}
\usepackage{geometry}                		% See geometry.pdf to learn the layout options. There are lots.
\geometry{letterpaper}                   		% ... or a4paper or a5paper or ... 
%\geometry{landscape}                		% Activate for for rotated page geometry
\usepackage[pdftex]{graphicx}				% Use pdf, png, jpg, or eps with pdflatex; use eps in DVI mode	
\usepackage{setspace}
\usepackage{physics}
\usepackage{subcaption}
\usepackage[numbers]{natbib}
\usepackage{pdfpages}
\usepackage{bm}
\usepackage{wrapfig}				% enables the use of \wrapfig, for figures with text wrapped around them
\usepackage{afterpage}
\newcommand\blankpage{%
    \null
    \thispagestyle{empty}%
    \addtocounter{page}{-1}%
    \newpage}

\usepackage{bibentry} % for list of publications
\usepackage{hyperref}

% \usepackage[english]{babel}
\graphicspath{{images/}{{\subfix{images/}}}}
\usepackage{tikz}
\usepackage{tikz-3dplot}
\usetikzlibrary{shapes,calc,positioning}
\tdplotsetmaincoords{70}{120}
\usepackage{siunitx}
\usepackage{adjustbox}
\usepackage{xcolor} %% Text color
\usepackage{blindtext}
\usepackage{subfiles} % Best loaded last in the preamble
\usepackage{multicol,multirow}  %% Figure in two columns
\usepackage{mathtools, nccmath}

\usepackage{lipsum}			% gives access to \lipsum, which dumps some latin text into your document as filler if you want to check formatting

%\usepackage[parfill]{parskip}    		% Activate to begin paragraphs with an empty line rather than an indent



% Here we set the page dimensions to match the standard thesis format.  These values should not be changed.
%%% SET LENTGH AND WIDTH %%%
\setlength{\textwidth}{6.5in}
\setlength{\textheight}{8.5in}
\setlength{\oddsidemargin}{0pt}
\setlength{\evensidemargin}{0pt}
\setlength{\topmargin}{0pt}
\setlength{\marginparsep}{0pt}
\setlength{\marginparwidth}{1in}


%\begin{document} starts LaTeX looking for actual content.  Everything above this point is purely formatting.
\begin{document}

\begin{titlepage}
\begin{center}

% \vspace* creates some vertical white space on the page to make the title page look more pleasing.  \vspace would do much the same thing, but would not insert the white space if we were at the top of a fresh page.  As this is the start of the document we're obviously at the beginning of a page, so the asterisk is necessary to ensure we still put in two cm of white space.
\vspace*{2cm}

{\Large \textbf{A Great Catchy Thesis Title}} % \huge sets the font size.  Other options include things like \large, \Large, \small, \tiny, etc.

\vspace{5cm}
{Thesis submitted in accordance with the requirements of
the University of Liverpool for the degree of Doctor in Philosophy\\ by}

\vspace{2cm}
{\large \textbf{H Diego Estrada-Lugo}\\ (Maximum Commander)}
% \vfill creates an arbitrary amount of vertical white space as necessary to fill the page
\vfill

April 2022

% \vspace*{3cm}

\end{center}
\end{titlepage}

% \frontmatter defines the pieces of the thesis which will use roman numerals for page numbering
\frontmatter 


% \chapter{} and/or \chapter*{} will create a chapter in your thesis.  Including the asterisk will cause the chapter to not appear in the table of contents.


\afterpage{\blankpage}

\chapter{Dedication}
\input{Chapters/Dedication}
\afterpage{\blankpage}

\chapter{Acknowledgements}
\input{Chapters/Acknowledgments}
\afterpage{\blankpage}

\chapter{Abstract}
\input{Chapters/Abstract}
\afterpage{\blankpage}

\chapter{Declaration}
\input{Chapters/Declaration}
\afterpage{\blankpage}

\chapter{List of Publications}
\input{Chapters/Publications}
\afterpage{\blankpage}




%\tableofcontents will create a table of contents.  By default it will include entries for any \chapter, \section, and \subsection command that appears in your thesis unless you have called the tag with an asterisk
\tableofcontents
\listoftables
\listoffigures

\chapter{Abbreviations and nomenclature}
\input{Chapters/nomenclature}


% The \mainmatter command defines the main body of the thesis and indicates where regular numbering starts
\mainmatter

\chapter{Introduction}
\input{Chapters/Introduction}

\chapter{Content chapter}
\input{Chapters/Chapter1}


\chapter{Conclusions}
\input{Chapters/Conclusions}

% the \appendix tag tells LaTeX where it should start labelling chapters with letters (denoting appendices) rather than numbers (denoting main chapters)
\appendix 


\chapter{Appendix title (if needed)}
\input{Chapters/Appendix}

\chapter{}
\input{Chapters/Appendix_two}



% \bibliographystyle command to choose the format of your bibliography. More examples of bibliography styles can be found at https://www.overleaf.com/learn/latex/Bibtex_bibliography_styles
\bibliographystyle{unsrtnat}

% \bibliography is the command for the actual file containing your bibliographic data. This file can be produced manually or automatically using software such as BibTeX. Both options can work, however, learning to use BibTex is beneficial in the long run. An example of the format needed to generate your own bibliography file can be found as the bibliography.bib file here provided.

\bibliography{bibliography}
\end{document}